% ===== PRÉAMBULE CANONIQUE — à coller VERBATIM en tête de CHAQUE .tex =====
\documentclass[11pt,a4paper]{article}
\usepackage[utf8]{inputenc}\usepackage[T1]{fontenc}\usepackage{lmodern}
\usepackage[french]{babel}
\usepackage{amsmath,amssymb,mathtools}\usepackage{icomma}
\usepackage{siunitx}\sisetup{locale=FR}  % \qty{}{}, \si{}, \ang{}
\usepackage{xcolor}\usepackage{colortbl}
\usepackage{tikz}\usepackage{pgfplots}\pgfplotsset{compat=1.18}
\usepackage[siunitx,europeanresistors]{circuitikz} % circuits (élec.)
\usepackage[most]{tcolorbox}\usepackage{enumitem}\usepackage{array,booktabs}
\usepackage[a4paper,margin=2.2cm]{geometry}
\usetikzlibrary{arrows.meta,calc,angles,quotes,positioning,patterns,%
  backgrounds,shapes.geometric,decorations.pathmorphing,decorations.markings,intersections}
\definecolor{primary}{RGB}{222,49,99}\definecolor{primarydark}{RGB}{150,22,55}
\definecolor{defcol}{RGB}{0,114,178}\definecolor{methcol}{RGB}{52,92,148}
\definecolor{excol}{RGB}{0,135,90}\definecolor{attncol}{RGB}{200,40,55}
\tcbset{boxbase/.style={enhanced,breakable,boxrule=0.9pt,arc=2.5pt,
  left=3mm,right=3mm,top=2.3mm,bottom=2.3mm,fonttitle=\bfseries,coltitle=white}}
% --- boîtes TUTORIEL (noms EXACTS, ne pas renommer) ---
\newtcolorbox{cle}[1][]{boxbase,colback=defcol!6,colframe=defcol,
  title={\textsc{À retenir}\ifx\relax#1\relax\relax\else\textmd{ : \itshape #1}\fi}}
\newtcolorbox{methode}[1][]{boxbase,colback=methcol!7,colframe=methcol,sharp corners,
  title={\textsc{Méthode}\ifx\relax#1\relax\relax\else\textmd{ --- \itshape #1}\fi}}
\newtcolorbox{exemplebox}[1][]{boxbase,colback=excol!5,colframe=excol,
  title={\textsc{Exemple}\ifx\relax#1\relax\relax\else\textmd{ : \itshape #1}\fi}}
\newtcolorbox{piege}[1][]{boxbase,colback=attncol!6,colframe=attncol,
  title={\textsc{Piège}\ifx\relax#1\relax\relax\else\textmd{ : \itshape #1}\fi}}
\newtcolorbox{remarque}[1][]{boxbase,colback=black!4,colframe=black!45,
  title={\textsc{Remarque}\ifx\relax#1\relax\relax\else\textmd{ : \itshape #1}\fi}}
% --- boîtes EXERCICE / CORRIGÉ (noms EXACTS, auto-comptées) ---
\newtcolorbox[auto counter]{exo}[1][]{boxbase,colback=primary!4,colframe=primary,
  title={\textsc{Exercice}~\thetcbcounter\ifx\relax#1\relax\relax\else\textmd{ --- \itshape #1}\fi}}
\newtcolorbox[auto counter]{corrige}[1][]{boxbase,colback=excol!4,colframe=excol,
  title={\textsc{Corrigé}~\thetcbcounter\ifx\relax#1\relax\relax\else\textmd{ --- \itshape #1}\fi}}
\newcommand{\vv}[1]{\vec{#1}}\newcommand{\ds}{\displaystyle}
\newcommand{\R}{\mathbb{R}}\newcommand{\N}{\mathbb{N}}\newcommand{\Z}{\mathbb{Z}}
% =========================================================================
\begin{document}
\usetikzlibrary{babel}
\begin{exo}[réseau mixte : série et parallèle]
Dans le réseau ci-dessous, $R_1=\qty{4}{\ohm}$ est en série avec l'association parallèle de
$R_2=\qty{12}{\ohm}$ et $R_3=\qty{12}{\ohm}$. Calcule la résistance équivalente entre les bornes.
\begin{center}
\begin{circuitikz}[scale=0.85,transform shape]
  \coordinate (A) at (2.3,0); \coordinate (B) at (4.6,0);
  \draw (0,0) to[short,o-] (0.5,0) to[R,l=$R_1$] (A);
  \draw (A) -- (2.3,1) to[R,l=$R_2$] (4.6,1) -- (B);
  \draw (A) -- (2.3,-1) to[R,l=$R_3$] (4.6,-1) -- (B);
  \draw (B) to[short,-o] (5.2,0);
  \fill (A) circle(1.4pt); \fill (B) circle(1.4pt);
\end{circuitikz}
\end{center}
\end{exo}

\begin{exo}[diviseur de courant]
Un courant $I=\qty{5}{\ampere}$ arrive sur deux résistances en parallèle, $\qty{4}{\ohm}$ et
$\qty{1}{\ohm}$.
\begin{center}
\begin{circuitikz}[scale=0.85,transform shape,>=Stealth]
  \coordinate (A) at (0,0); \coordinate (B) at (3.6,0);
  \draw[thick] (A)--(0,1); \draw (0,1) to[R,l=$\qty{4}{\ohm}$] (3.6,1); \draw[thick] (3.6,1)--(B);
  \draw[thick] (A)--(0,-1); \draw (0,-1) to[R,l=$\qty{1}{\ohm}$] (3.6,-1); \draw[thick] (3.6,-1)--(B);
  \draw[->,primary,line width=1pt] (-1.5,0)--(-0.12,0) node[midway,above,font=\footnotesize,primary]{$I$};
  \draw (B) -- (4.2,0);
  \fill (A) circle(1.4pt); \fill (B) circle(1.4pt);
\end{circuitikz}
\end{center}
\begin{enumerate}[label=\alph*)]
  \item Les deux résistances ont la même tension. Calcule-la (via $R_{\text{eq}}$).
  \item En déduire l'intensité dans chaque branche. Laquelle « prend » le plus de courant ?
\end{enumerate}
\end{exo}

\begin{exo}[générateur réel débitant dans une résistance]
Une pile de f.é.m. $E=\qty{12}{\volt}$ et de résistance interne $r=\qty{0.5}{\ohm}$ alimente une
résistance $R=\qty{5.5}{\ohm}$.
\begin{center}
\begin{circuitikz}[scale=0.8,transform shape]
  \draw (0,0) to[battery1,l_=$E$] (0,2) to[R,l_=$r$] (2,2) -- (4,2) to[R,l=$R$] (4,0) -- (0,0);
  \node[primarydark,font=\scriptsize] at (1,2.6){générateur réel};
\end{circuitikz}
\end{center}
\begin{enumerate}[label=\alph*)]
  \item Calcule l'intensité $I$ débitée.
  \item En déduire la tension $U$ aux bornes de la pile.
\end{enumerate}
\end{exo}

\begin{exo}[court-circuit d'une pile]
Une pile de f.é.m. $E=\qty{4.5}{\volt}$ a une résistance interne $r=\qty{0.5}{\ohm}$. On la met en
\textbf{court-circuit} (on relie ses bornes par un fil de résistance quasi nulle).
\begin{enumerate}[label=\alph*)]
  \item Calcule l'intensité de court-circuit $I_{\text{cc}}$.
  \item Calcule la puissance dissipée dans la résistance interne. Que risque-t-il de se passer ?
\end{enumerate}
\end{exo}

\begin{exo}[trois résistances en parallèle : intensité totale]
Trois résistances $R_1=\qty{10}{\ohm}$, $R_2=\qty{100}{\ohm}$ et $R_3=\qty{20}{\ohm}$ sont en
parallèle sous une tension $U_{AB}=\qty{2.5}{\volt}$. Calcule l'intensité totale $I$ de deux façons :
via chaque branche, puis via $R_{\text{eq}}$.
\end{exo}

\end{document}
